%% ch01-architecture.tex — System Architecture
%% JA4proxy Technical Reference Manual

\chapter{System Architecture}
\index{architecture}
\label{ch:architecture}

\begin{chapterquote}
\ja\ makes allow/block/tarpit decisions based entirely on plaintext metadata visible
before and during the TLS handshake. It never decrypts traffic, never holds TLS keys,
and forwards allowed connections byte-for-byte unchanged.
\end{chapterquote}

\newthought{Three design principles} govern every architectural decision in \ja:
\textbf{fail open} (a missed threat is recoverable; a blocked legitimate user is not),
\textbf{never decrypt} (the proxy is cryptographically transparent — it reads only the
plaintext ClientHello), and \textbf{monitor first} (the default dial is 0; the system
observes and scores before it ever blocks).

\section{Design Principles}
\index{design principles}
\index{fail open}

\subsection{The Asymmetry Principle}
\index{asymmetry principle}
\index{false positive}
\index{false negative}

The most important design constraint in \ja\ is the cost asymmetry between the two
error types:

\begin{fullwidth}
\begin{table}[h]
\caption{Error cost asymmetry — the governing principle of all threshold and TTL decisions}
\label{tab:asymmetry}
\begin{tabular}{@{}lllp{7cm}@{}}
\toprule
\textbf{Error type} & \textbf{Example} & \textbf{Cost} & \textbf{Consequence} \\
\midrule
False negative & Bad bot slips through during cache sync window & \textcolor{ja4green}{Low} & Recoverable. Log entry exists; retrospective analysis possible; no user harm. \\
\addlinespace
False positive & Real browser is blocked or tarpitted & \textcolor{ja4red}{\textbf{High}} & Not recoverable. User leaves. Revenue impact. Reputation damage. \\
\bottomrule
\end{tabular}
\end{table}
\end{fullwidth}

Every caching TTL, every threshold default, every fallback behaviour, and every new
feature decision flows from this asymmetry. When the answer is not obvious, fail open.

\begin{importantbox}[The Asymmetry Rule]
ALLOW decisions are cached for 30 minutes. BLOCK decisions are cached for 30 seconds.
When Redis says block but the local in-process cache says allow, the local cache wins.
If Redis is down, real browsers continue to work.
\end{importantbox}

\subsection{Never Decrypt}
\index{TLS!never decrypt}
\index{ClientHello}

\ja\ operates entirely on \textbf{plaintext metadata}. The TLS ClientHello — the
first message a TLS client sends — is transmitted in plaintext before any encryption
is established. It contains:

\begin{itemize}
  \item The TLS version and list of supported cipher suites
  \item Supported extensions and their parameters
  \item The SNI (Server Name Indication) hostname
  \item ALPN (Application-Layer Protocol Negotiation) values
  \item Elliptic curve groups and point formats
\end{itemize}

From these fields, the JA4 fingerprinting algorithm produces a deterministic string
that identifies the TLS client implementation (browser, bot, scanner, malware C2)
with high accuracy. The proxy never participates in the TLS handshake — it is a
passthrough observer and gatekeeper, not a terminator.

\subsection{Monitor First}
\index{monitor mode}
\index{dial!default}

The default configuration ships with \configkey{dial: 0}. At dial=0, the pipeline
runs completely — all signals are collected, all scores are computed, all decisions
are made — but no connection is ever blocked. Every connection is allowed with a
\badgemonitor\ annotation in the log.

This means operators can deploy \ja\ in front of production traffic immediately,
observe the score distribution, tune thresholds, and build confidence before enabling
any blocking. Raising the dial is a conscious operator decision.

\subsection{Config-Driven, Hot-Reloadable}
\index{hot reload}
\index{SIGHUP}

Every behavioural parameter is configurable in \code{config/proxy.yml}. The proxy
watches for \code{SIGHUP} and reloads configuration without dropping connections.
The Management UI triggers reload via a Redis pub/sub message. After reload, new
configuration applies to all subsequent connections.

Only the listen port, Redis URL, and TLS certificate paths require a process restart.
All other parameters — dial, bypass conditions, thresholds, blacklists, whitelists —
are hot-reloadable.\sidenote{Config reload events are logged and reflected in the
\code{ja4proxy\_config\_reloads\_total} Prometheus counter.}

\section{Component Overview}
\index{components}

\begin{fullwidth}
\begin{table}[h]
\caption{System components and their responsibilities}
\label{tab:components}
\begin{tabular}{@{}lllp{6.5cm}@{}}
\toprule
\textbf{Component} & \textbf{Language} & \textbf{Port} & \textbf{Responsibility} \\
\midrule
HAProxy & C & :443 & TLS passthrough load balancer; PROXY protocol v2 injection; upstream TLS listener \\
\addlinespace
\ja\ Proxy & Python / Go\textsuperscript{*} & :8080 & TLS ClientHello sniffing; pipeline execution; allow/block/tarpit/ban decisions; Redis I/O \\
\addlinespace
Redis & Redis 7 & :6379 & Shared state: bans, rate windows, beaconing timestamps, JA4 lists, HLL per CIDR, event stream \\
\addlinespace
Analytics Node & Python & — & Reads Redis Streams (\code{XREADGROUP}); cross-instance aggregation; campaign detection; writes findings to Redis \\
\addlinespace
Prometheus & Go & :9090 & Metrics collection and storage \\
\addlinespace
Grafana & Go & :3001 & Dashboards; visualisation of proxy telemetry \\
\addlinespace
Loki & Go & :3100 & Structured log aggregation \\
\addlinespace
Alertmanager & Go & :9093 & Alert routing and deduplication \\
\addlinespace
Backend & any & :443 & Protected origin service; receives only allowed, unmodified TLS connections \\
\bottomrule
\end{tabular}
{\footnotesize\textsuperscript{*}Python is the current production path.
The Go rewrite (\code{cmd/proxy/}) reached feature parity in Phase~15.}
\end{table}
\end{fullwidth}

\section{Network Topology}
\index{network topology}
\index{network zones}
\index{DMZ}

\subsection{Traffic Flow Diagram}

\begin{figure}[h]
\begin{fullwidth}
\begin{tikzpicture}[
  node distance=1.6cm,
  box/.style={draw=ja4navy, fill=ja4lightgray, rounded corners=3pt,
              minimum width=2.6cm, minimum height=0.9cm,
              font=\small\sffamily, text centered},
  internet/.style={draw=ja4red, fill=ja4redlight, rounded corners=3pt,
                   minimum width=2.6cm, minimum height=0.9cm,
                   font=\small\sffamily, text centered},
  backend/.style={draw=ja4green, fill=ja4greenlight, rounded corners=3pt,
                  minimum width=2.6cm, minimum height=0.9cm,
                  font=\small\sffamily, text centered},
  redis/.style={draw=ja4purple, fill=ja4purplelight, rounded corners=3pt,
                minimum width=2.2cm, minimum height=0.8cm,
                font=\small\sffamily, text centered},
  arrow/.style={->, >=Stealth, thick, draw=ja4navy},
  label/.style={font=\footnotesize\sffamily\itshape, text=ja4gray},
]
  \node[internet] (internet) {Internet};
  \node[box, right=of internet] (haproxy) {HAProxy\\{\footnotesize :443}};
  \node[box, right=of haproxy] (proxy1) {\ja\ \#1\\{\footnotesize :8080}};
  \node[box, below=0.5cm of proxy1] (proxy2) {\ja\ \#2\\{\footnotesize :8080}};
  \node[backend, right=2.2cm of proxy1] (backend) {Backend\\{\footnotesize :443}};
  \node[redis, below=1.4cm of haproxy] (redis) {Redis\\{\footnotesize :6379}};
  \node[box, right=of redis] (analytics) {Analytics\\{\footnotesize Node}};

  \draw[arrow] (internet) -- node[above, label]{TLS} (haproxy);
  \draw[arrow] (haproxy) -- node[above, label]{TCP} (proxy1);
  \draw[arrow] (haproxy) -- (proxy2);
  \draw[arrow] (proxy1) -- node[above, label]{TLS} (backend);
  \draw[arrow] (proxy2) -- (backend);
  \draw[arrow, <->] (proxy1) -- node[right, label]{state} (redis);
  \draw[arrow, <->] (proxy2) -- (redis);
  \draw[arrow] (redis) -- node[above, label]{Streams} (analytics);
  \draw[arrow] (analytics) -- node[below, label]{findings} (redis);
\end{tikzpicture}
\end{fullwidth}
\caption{High-level traffic flow. HAProxy injects PROXY protocol v2 headers so \ja\ can recover the real client IP.}
\label{fig:traffic-flow}
\end{figure}

\subsection{Network Zones}
\index{network zones!DMZ}
\index{network zones!App}
\index{network zones!Data}
\index{network zones!Management}

A production deployment uses four network zones:

\begin{fullwidth}
\begin{table}[h]
\caption{Network zone definitions and permitted traffic}
\label{tab:zones}
\begin{tabular}{@{}llp{4cm}p{5cm}@{}}
\toprule
\textbf{Zone} & \textbf{Subnet (example)} & \textbf{Contains} & \textbf{Permitted inbound} \\
\midrule
DMZ & 10.0.1.0/24 & HAProxy load balancer & Internet :443 only \\
App & 10.0.2.0/24 & \ja\ proxy instances & HAProxy TCP from DMZ only \\
Data & 10.0.3.0/24 & Redis, analytics node & App zone on :6379 only \\
Management & 10.0.4.0/24 & Prometheus, Grafana, Loki & Operator VPN only \\
\bottomrule
\end{tabular}
\end{table}
\end{fullwidth}

The Backend origin server sits behind the App zone, reachable only from \ja\ instances.
Internet traffic never reaches the App zone directly. The Data zone is not routable from
the DMZ or from the internet.

\section{Data Flow: TLS Connection Lifecycle}
\index{TLS!connection lifecycle}
\index{ClientHello!parsing}

\subsection{Stage 1 — TCP Accept and IP Extraction}

When HAProxy accepts a TLS connection, it wraps it in a PROXY protocol v2
header\sidenote{PROXY protocol v2 is a binary framing format prepended by HAProxy to
the TCP stream. It carries the original client IP:port and the HAProxy IP:port.} and
forwards the raw TCP stream to a \ja\ instance on port 8080.

\ja\ reads the PROXY protocol header first, extracting:
\begin{itemize}
  \item The \textbf{real client IP} (IPv4 or IPv6)
  \item The real client port
  \item The HAProxy (sender) IP — used for upstream CIDR trust verification
\end{itemize}

The client IP is canonicalised to its compressed string representation
(\code{ipaddress.ip\_address(ip).compressed} in Python, \code{netip.Addr.String()} in Go)
immediately. All subsequent pipeline steps use this canonical form.

\subsection{Stage 2 — TLS ClientHello Sniff}
\index{TLS!ClientHello!sniff}

After PROXY protocol extraction, \ja\ reads the first TLS record from the TCP stream.
This is always the ClientHello — a plaintext record (record type 0x16, content type 0x01).
The proxy parses:

\begin{itemize}
  \item Handshake type and TLS record version
  \item Client random (32 bytes, not used for fingerprinting)
  \item Session ID length and value
  \item Cipher suite list (ordered)
  \item Compression methods
  \item Extensions: SNI (0x0000), ALPN (0x0010), supported groups (0x000a),
        EC point formats (0x000b), signature algorithms (0x000d),
        supported versions (0x002b), PSK key exchange modes (0x002d),
        key share (0x0033), and others
\end{itemize}

From this data, the JA4, JA4X, and JA4T fingerprints are computed (see
Chapter~\ref{ch:signals}). The entire sniff operation is non-blocking and completes
before any response is sent.

\subsection{Stage 3 — Pipeline Execution}

The pipeline (Chapter~\ref{ch:pipeline}) runs all 12 layers. Layers 0–2 are
bypass/block decisions that short-circuit the pipeline immediately. Layers 4–8 are
signals that run in parallel via \code{asyncio.gather()}. The pipeline produces a
score (0–100) and an action.

\subsection{Stage 4 — Action Execution}

Based on the action from the pipeline:

\begin{description}
  \item[\badgeallow] The buffered ClientHello bytes are forwarded to the backend. \ja\
    acts as a transparent TCP relay for the remainder of the connection.
  \item[\badgeblock] A TCP RST is sent. The connection is closed immediately.
  \item[\badgetarpit] The connection is accepted but data is delivered to the backend
    at an artificially low rate (or deliberately delayed), consuming attacker resources.
  \item[\badgeban] The IP is added to Redis with a TTL. Subsequent connections from
    this IP are rejected immediately at layer 0 without pipeline execution.
  \item[Flag] The connection is allowed but a high-priority alert is emitted to the
    Redis Stream for the Analytics node.
\end{description}

\subsection{Stage 5 — Event Emission}

Every connection that completes pipeline evaluation produces a JSON event written to
the \code{events} Redis Stream via \code{XADD} (fire-and-forget). The analytics node
consumes these events for cross-instance aggregation, campaign detection, and score
drift alerting. Stream writes add zero latency to the hot path.

\section{IPv6 Handling}
\index{IPv6}
\index{IPv4}

IPv6 is a first-class concern throughout \ja. Every feature that touches an IP address
handles both address families:

\begin{fullwidth}
\begin{table}[h]
\caption{IPv6 handling rules by subsystem}
\label{tab:ipv6}
\begin{tabular}{@{}lll@{}}
\toprule
\textbf{Subsystem} & \textbf{IPv4 rule} & \textbf{IPv6 rule} \\
\midrule
Rate limiting keys & Full IP string & Full IP string (no truncation) \\
CIDR matching & In-process pytricia / Go trie & Same trie handles both natively \\
GeoIP lookup & MaxMind GeoLite2 City/ASN & Same DB covers IPv6 \\
HyperLogLog per subnet & \code{hll:cidr24:\{cidr\}} (/24) & \code{hll:cidr48:\{cidr\}} (/48) \\
Beaconing key & \code{beacon:\{ip\}:\{ja4\}} & Same key format, full IPv6 \\
Ban key & \code{ban:\{ip\}} & Same key format, full IPv6 \\
RDAP lookup & RIR via IANA bootstrap & Same bootstrap — works for IPv6 \\
Block expansion & Never beyond /24 & Never beyond /48 \\
\bottomrule
\end{tabular}
\end{table}
\end{fullwidth}

\begin{notebox}[IPv6 /48 Equivalence]
A /48 IPv6 prefix is the approximate equivalent of an IPv4 /24 for population density
purposes. A typical ISP assigns a /48 to a single customer premises. Block expansion
(RDAP Phase~11) never expands beyond /48 for IPv6 or /24 for IPv4 — expanding further
would risk blocking entire ISP customer populations.
\end{notebox}

\section{The Python and Go Implementations}
\index{Go rewrite}
\index{Python implementation}

\subsection{Python (Current Production Path)}

The Python implementation in \code{proxy.py} and \code{src/security/} is the current
production-validated path. It uses \code{asyncio} throughout for non-blocking I/O.
All external service calls use \code{asyncio.create\_task()} — fire-and-forget from
the hot path. No blocking I/O is permitted on the connection handling code path.

\subsection{Go Rewrite (Phase 15)}
\index{Go rewrite!Phase 15}

The Go rewrite in \code{cmd/proxy/} targets 10–50× throughput improvement over the
Python implementation. It has achieved feature parity with the Python implementation
including JA4X fingerprinting and Lua policy embedding. The Go implementation uses
\code{net/netip} for IP handling, Go's \code{crypto/tls} stack for ClientHello
parsing, and goroutine-per-connection concurrency.

\begin{notebox}[Implementation Status]
Python is \statuspoc\ — functional and tested. Go rewrite is \statuspoc\ — feature
parity achieved but not yet production-hardened. Neither implementation has completed
the full Phase~14 production hardening checklist.
\end{notebox}

\section{Horizontal Scaling}
\index{horizontal scaling}
\index{Redis!shared state}

Multiple \ja\ instances share state through Redis. Because all ban decisions, rate
windows, JA4 lists, and beaconing data live in Redis, any instance can enforce
decisions made by any other instance. There is no sticky routing requirement — HAProxy
distributes connections across instances with standard round-robin or least-connection
load balancing.

The only per-instance state is the in-process LRU cache (\code{src/cache/local\_cache.py}).
This cache is intentionally not synchronised across instances. Cache misses are handled
by Redis fallback. The local cache uses conservative TTLs (30min for ALLOW, 30s for
BLOCK) to bound the window in which a newly added ban is not yet enforced by a
particular instance.
